(1) 理想金属边界要求壁面切向电场为零，TE 模可由纵向磁场 \(H_z\propto\cos(m\pi x/a)\cos(n\pi y/b)\) 描述，整数 \(m,n\ge0\) 不能同时为零。横向本征波数 \(k_c^2=(m\pi/a)^2+(n\pi/b)^2\)；轴向传播常数 \(\beta\) 与频率满足 \(\beta^2=\omega^2\mu\varepsilon-k_c^2=\omega^2/c^2-k_c^2\)。
(2) 当 \(a>b\) 时最小非零横向本征值为 \(k_c=\pi/a\)，对应 TE10 模；因此 \(\omega_c=\pi c/a\)，普通频率 \(f_c=c/(2a)\)。若 \(\omega<\omega_c\)，\(\beta\) 为纯虚数，轴向场倏逝，无法把它解释为远距离携能传播的普通行波。
(3) 在 \(\omega>\omega_c\)，\(\beta=(1/c)\sqrt{\omega^2-\omega_c^2}\)，导波波长 \(\lambda_g=2\pi/\beta\)。相速度 \(v_p=\omega/\beta=c[1-(\omega_c/\omega)^2]^{-1/2}\)，群速度 \(v_g=(d\beta/d\omega)^{-1}=c\sqrt{1-(\omega_c/\omega)^2}\)，乘积 \(v_pv_g=c^2\)，与无损波导的能量传播速度相符。
(4) 频率趋近截止上方时 \(\beta\to0\)、\(\lambda_g\to\infty\)、\(v_p\to\infty\) 而 \(v_g\to0\)；高于截止很多时 \(v_p,v_g\to c\)。超大的相速度反映等相位面沿波导的几何传播，并不对应信息或能量传递速率；实际脉冲的传播由群速度以及完整色散与因果响应控制。
